% ID: SEL-D0049-OB-E02
% ESTADO: validada. S050.
\textbf{a)} $f'(x)=x^2+x-2=(x+2)(x-1)$. La derivada es positiva fuera de $[-2,1]$ y negativa dentro. El signo cambia de positivo a negativo en $-2$, máximo relativo de valor $19/3$, y de negativo a positivo en $1$, mínimo relativo de valor $11/6$.

\textbf{b)} Para $g=f'$, el punto es $(2,g(2))=(2,4)$ y la pendiente es $g'(2)=2\cdot2+1=5$. La tangente es $\boxed{y=5x-6}$.

\textbf{c)} La parábola $g$ abre hacia arriba, corta el eje horizontal en $-2$ y $1$ y tiene vértice $(-1/2,-9/4)$. La recta pasa por $(2,4)$ y queda por debajo de la parábola porque $g(x)-(5x-6)=(x-2)^2$.
\begin{center}
\setlength{\unitlength}{0.8mm}
\begin{picture}(98,65)(-7,-5)
\qbezier(0,12)(42,12)(84,12)
\qbezier(36,0)(36,28)(36,56)
\put(85,11){$x$}
\put(38,55){$y$}
\qbezier(0,28)(36,-32)(72,52)
\qbezier(48,8)(60,28)(72,48)
\put(3,29){\scriptsize $g$}
\put(73,45){\scriptsize $y=5x-6$}
\put(60,24){\scriptsize $(2,4)$}
\put(10,7){\scriptsize $-2$}
\put(47,14){\scriptsize $1$}
\end{picture}
\end{center}
